\documentclass[12pt,letterpaper]{article}
\usepackage[margin=1in]{geometry}
\usepackage{amsmath,amssymb,amsthm,blkarray,braket,caption,enumitem,fancyvrb,graphicx,hyperref,relsize,verbatim}
\usepackage[numbers]{natbib}
\usepackage[section]{placeins}
\graphicspath{{./images/}}
\bibliographystyle{apalike}

\begin{document}

\title{An Introduction to Measure Theory - Chapter 1}
\author{Emil Marinov}
\maketitle

\section{Exercise 1.1.1 - Boolean Closure}

\textit{Show that if $E, F \subset \mathbb{R}^d$ are elementary sets, then the union $E \cup F$, 
the intersection $E \cap F$, and the
set theoretic difference $E \backslash F := \{x \in E : x \notin F\}$, and the symmetric
difference $E \Delta F := (E \backslash F) \cup (F \backslash E)$ are also elementary. If $x \in \mathbb{R}^d$, 
show that the translate $E + x := \{y+x : y \in E\}$ is also an elementary
set.}
\\
\\
An \textbf{elementary set} is that which is a union of a finite amount of boxes.  Thus, let
$\bigcup_i E_i$ and $\bigcup_j F_j$ be the unions of the finite number of boxes comprising sets $E$ and $F$, respectively.

\begin{equation}
   \begin{aligned}
        E \cup F = \bigcup_i E_i \cup \bigcup_j F_j
        = \bigcup_i \bigcup_j E_i \cup F_j
   \end{aligned}
\end{equation}

\begin{equation}
   \begin{aligned}
      E \cap F = \bigcup_i E_i \cap \bigcup_j F_j
      = \bigcup_i \bigcup_j E_i \cap F_j
   \end{aligned}
\end{equation}

\begin{equation}
   \begin{aligned}
      E \backslash F = \bigcup_i E_i \backslash \bigcup_j F_j 
      = \bigcup_i \left( E_i \cap \neg \bigcup_j F_j \right)
      = \bigcup_i \left( E_i \backslash \bigcup_j F_j \right)
   \end{aligned}
\end{equation}

This one tripped me up a bit because elementary sets aren't closed under set negation,
and I was inherently proving the closure of elementary sets under the set difference for some reason.

\begin{equation}
   \begin{aligned}
      E \Delta F = E \backslash F \cup F \backslash E
      = \bigcup_i \left( E_i \backslash \bigcup_k F_k \right) \cup \bigcup_j \left( F_j \backslash \bigcup_k E_k \right) 
      = \bigcup_i \bigcup_j \left( E_i \backslash \bigcup_k F_k \cup F_j \backslash \bigcup_k E_k \right)
   \end{aligned}
\end{equation}

\begin{equation}
   \begin{aligned}
      E + x = \bigcup_i \left( E_i + x \right)
   \end{aligned}
\end{equation}

The set theoretical operations of union, intersection, difference, and symmetric difference, and translation, 
between two elementary sets yield a union of finitely countable boxes, which is itself an elementary set. $\square$

\section{Exercise 1.1.2}

\textit{Give an alternate proof of Lemma 1.1.2(ii) by showing that any two partitions of $E$ into boxes admit a mutual refinement
into boxes that arise from taking Cartesian products of elements from
finite collections of disjoint intervals.}
\\
\textbf{Lemma 1.1.1(ii)}
If $E$ is partitioned as the finite union $B_1 \cup ... \cup B_k$ of disjoint
boxes, then the quantity $m(E) := |B_1| + ... + |B_k|$ is independent of the partition. In other words, given any other
partition $B'_1 \cup ... \cup B'_{k'}$ of $E$, one has $|B_1| + ... + |B_k| = |B'_1| + ... + |B'_{k'}|$.
\\
\begin{equation}
   \begin{aligned}
      \lvert E \rvert = \lvert \bigcup_k B_k \rvert = \lvert \bigcup_{\tilde{k}} B_{\tilde{k}} \rvert \\
      = \lvert \bigcup_k I_{k,1} \times ... \times I_{k,d} \rvert = \lvert \bigcup_{\tilde{k}} I_{\tilde{k},1} \times ... \times I_{\tilde{k},d} \rvert
   \end{aligned}
\end{equation}

Let $\{ J_{m,1}, ..., J_{m,d} \}$ be the disjoint intervals between all $2k$ endpoints of $\{ I_{k,1}, ..., I_{k,d} \}$ and 
$\{ \tilde{J}_{\tilde{m},1}, ..., \tilde{J}_{\tilde{m},d} \}$ be the disjoint intervals between all $2\tilde{k}$
endpoints of $\{ I_{\tilde{k},1}, ..., I_{\tilde{k},d} \}$.
Let $l_{k, i}$ for $i \in [d]$ denote the disjoint intervals from $J_{m,i}$ that union to form $I_{k,i}$
and $\tilde{l}_{\tilde{k}, i}$ denote the disjoint intervals from $\tilde{J}_{\tilde{m},i}$ that union to form $I_{\tilde{k},i}$.

\begin{equation}
   \begin{aligned}
      \lvert E \rvert
      = \lvert \bigcup_k \left( \bigcup_{l_{k,1}} J_{l_{k,1}} \times ... \times \bigcup_{l_{k,d}} J_{l_{k,d}} \right) \rvert = 
      \lvert \bigcup_{\tilde{k}} \left( \bigcup_{\tilde{l}_{\tilde{k},1}} J_{\tilde{l}_{\tilde{k}, 1}} \times ... \times \bigcup_{\tilde{l}_{\tilde{k},d}} J_{\tilde{l}_{\tilde{k}, d}} \right) \rvert
   \end{aligned}
\end{equation}

Thus, it has been shown that any two partitions of $E$ into boxes admit a mutual refinement into boxes that arise from taking the 
Cartesian product from elements from finite collections of disjoint intervals. $\square$

\section{Exercise 1.1.2.5}

\textit{Give some examples of sets for which the limit formula $m(e) = \frac{1}{N} \#\left( E \cap \frac{1}{N} \mathcal{N}^d \right)$ does not exist.}

\textbf{Example 1}.  The limit can fail if it diverges to $\pm \infty$.
This occurs when the interval is unbounded.

\begin{equation}
   \begin{aligned}
      E = \mathbb{R}^d
   \end{aligned}
\end{equation}

\textbf{Example 2}.  Even if bounded, the limit can fail if it oscillates.

\begin{equation}
   \begin{aligned}
      E = \bigcup_{k=1}^\infty \left( \frac{1}{2^{2k}}, \frac{1}{2^{2k-1}} \right)
   \end{aligned}
\end{equation}

This is the set of alternating dyadic intervals close to 0.

\section{Exercise 1.1.3 - Uniqueness of Elementary Measure}

\textit{Let $d \ge 1$.  Let $m' : \mathcal{E}(\mathbb{R}^d) \to \mathbb{R}^+$ be a map from the collection 
$\mathcal{E}(\mathbb{R}^d)$ of elementary subsets of $\mathbb{R}^d$ to the nonnegative reals that obeys the non-negativity,
finite additivity, and translation invariance properties. Show that
there exists a constant $c \in \mathbb{R}^+$ such that $m'(E) = cm(E)$ for all elementary sets $E$.
In particular, if we impose the additional normalisation $m'[0,1)^d = 1$ then $m' \equiv m$.
Hint: Set $c := m'([0,1)^d)$, and then compute $m'([0,\frac{1}{n})^d)$ for any positive integer $n$.}\\
\\

The half-open sets in the hint indicate that finite additivity corresponding to tiling is involved in this proof.
Further, the $\frac{1}{n}$ in the hint indicates that the set to be measured is broken into $n^d$ tiles.

Let $c = m'([0,1)^d)$ be a scalar corresponding to the $m'$ measure of the unit box.
Partition the unit box into $n^d$ tiles as: $[0,1) = \bigcup_{i=1}^{n^d} \left( [0,1/n)^d + x_i \right)$.
Then,

\begin{equation}
   \begin{aligned}
      m'([0,1)^d) = m'\left(\bigcup_{i=1}^{n^d} \left( [0,1/n)^d + x_i \right)\right) \\
      = \sum_{i=1}^{n^d} m'\left([0,1/n)^d + x_i \right) && \text{By finite additivity.} \\
      = \sum_{i=1}^{n^d} m'\left([0,1/n)^d \right) && \text{By translational invariance.} \\
      = n^d m'([0,1/n)^d)
   \end{aligned}
\end{equation}

Similarly, $m([0,1)^d) = n^d m([0,1/n)^d)$.
Rewrite $c$ as the ratio of these two measures:

\begin{equation}
   \begin{aligned}
      c = \frac{n^d m'([0,1/n)^d)}{n^d m([0,1/n)^d)} \\
      = \frac{m'([0,1/n)^d)}{m([0,1/n)^d} \\
      = \frac{m'([0,1)^d)}{m([0,1)^d)}
   \end{aligned}
\end{equation}

Thus, by finite addivitiy of the elementary set's decomposition into finite disjoint unions of half-open boxes,
the existence of a unique $c$ for all $E \in \mathcal{E}^d$ has been proved. $\square$

\subsubsection{Scratch Work}
From non-negativity and finite additivity, monotonicity is derived: 
Let $X$ be a subset of $Y$.  Then, there exists some elementary set $Z = Y \backslash X$ disjoint to $X$ such that $Y = X \cup Z$.
\begin{equation}
   \begin{aligned}
      m(Y) = m(X) + m(Z) \quad \forall X,Z : X \cap Z = \emptyset && \text{By finite addivity.} \\
      m(X) = m(Y) - m(Z) \\
      m(X) \le m(Y) && \text{By non-negativitity of $Z$.}
   \end{aligned}
\end{equation}

From this, finite subadditivity may be derived:

\begin{equation}
   \begin{aligned}
      m(X \cup Y) = m(X \cup (Y \backslash X)) \\ 
      = m(X) + m(Y \backslash X) && \text{By finite additivity.} \\
      \le m(X) + m(Y) && \text{By monotonicity of $m(Y \backslash X)$ and $m(Y)$.}
   \end{aligned}
\end{equation}

Now, let $E$ be the union of a finite amount of elementary sets, $E = \bigcup_i E_i$.
By finite subadditivity, $m(E) \le \sum_i E_i$, and by scalar multiplication, $c * m(E) \le c * \sum_i E_i$

Let $c = m'([0,1)^d)$.

\begin{equation}
   \begin{aligned}
      c m(E) = m'([0,1)^d) m(E) \\
      \le m'([0,1)^d) \sum_i m(E_i) \\
      = \sum_i m'([0,1)^d) m(E_i)
   \end{aligned}
\end{equation}

\begin{equation}
   \begin{aligned}
      m'([0, \frac{1}{n})^d) = m'([0,1)^d) m([0, \frac{1}{n})^d)
    \end{aligned}
\end{equation}

\end{document}